% Anexos (cargado por main.tex tras \appendix)

\chapter{Glosario de términos}
\label{anexo:glosario}

Glosario pensado para el lector no familiarizado con las criptomonedas ni con la
microestructura de mercado. Recoge los términos empleados a lo largo de la memoria.

\begin{description}[align=left, leftmargin=3.2cm, labelwidth=3cm, font=\bfseries]
\item[Mercado de predicción] Mercado en el que se negocian contratos cuyo valor depende del
resultado de un acontecimiento futuro.
\item[Contrato binario] Contrato que paga una unidad si el suceso ocurre y cero si no; su
precio vive entre 0 y 1 y se interpreta como probabilidad.
\item[Polymarket] Plataforma de mercados de predicción descentralizada, con emparejamiento
fuera de cadena y liquidación en cadena.
\item[Libro de órdenes (LOB)] Lista, ordenada por precio, de las órdenes de compra y venta
pendientes de ejecución.
\item[Bid / Ask] Mejor precio de compra (\emph{bid}) y mejor precio de venta (\emph{ask}).
\item[Spread] Diferencia entre \emph{ask} y \emph{bid}; primer coste de cruzar el mercado.
\item[Profundidad] Volumen disponible en cada nivel de precio del libro.
\item[Precio medio (\emph{mid})] Punto medio entre \emph{bid} y \emph{ask}.
\item[Tick] Variación mínima de precio del mercado; unidad de medida de los resultados.
\item[Microestructura] Estudio de cómo se forman los precios a partir de órdenes, \emph{spread},
profundidad y flujo, a escalas de tiempo muy finas.
\item[Markout] Movimiento del precio desde el instante de ejecución a lo largo de un horizonte;
sirve para medir si una operación fue buena.
\item[Latencia] Retardo entre observar la oportunidad y ejecutar la orden.
\item[Spot / Perpetuo] Mercado de Bitcoin al contado (\emph{spot}) y de futuros sin vencimiento
(\emph{perpetual futures}); referencias externas del contrato.
\item[Oráculo de precios] Fuente externa que aporta el precio de referencia de Bitcoin.
\item[Volatilidad realizada] Medida de cuánto se ha movido el precio recientemente.
\item[Ventana \emph{prestart}] Intervalo inmediatamente anterior a la apertura del mercado.
\item[Validación temporal] Partición de los datos respetando el orden del tiempo (entrenar con
el pasado, evaluar con el futuro).
\item[Fuera de muestra (OOS)] Datos no usados ni para entrenar ni para seleccionar.
\item[Fuera de pliegue (OOF)] Protocolo en el que las predicciones que consume una segunda etapa
proceden siempre de un modelo que no las entrenó.
\item[Fuga temporal] Filtración de información del futuro hacia el entrenamiento o las
características; principal fuente de resultados ilusorios.
\item[Congelación] Fijar el modelo y los umbrales antes de observar el conjunto fuera de muestra.
\item[Drawdown] Mayor caída acumulada desde un máximo previo.
\item[\emph{Bootstrap}] Remuestreo con reemplazo para estimar la incertidumbre de un resultado.
\item[Paper-shadow] Seguimiento en sombra: se registra lo que la política habría hecho, sin
arriesgar capital.
\item[Gradient boosting] Familia de modelos que suma muchos árboles de decisión, cada uno
corrigiendo los errores del anterior.
\end{description}

\chapter{Conjunto de características}
\label{anexo:features}

El especialista final emplea 36 características causales sin variables de reloj
(\emph{no\_clock}), agrupadas por bloques:

\begin{itemize}[noitemsep]
\item \textbf{Estado local del libro de Polymarket:} precio medio, \emph{spread}, microprice,
desbalance de órdenes, coste visible de entrada y deltas recientes del precio.
\item \textbf{Referencias externas:} retornos orientados del mercado al contado y de los
futuros perpetuos a 2 y 8 segundos; medidas de consenso y de actualización.
\item \textbf{Volatilidad:} volatilidad realizada del perpetuo a 5 segundos (también usada por
el filtro de volatilidad).
\item \textbf{Contexto de ventana:} fase de la ventana \emph{prestart} y proxies de calidad del
libro (cobertura, indicadores de degradación).
\end{itemize}

Se excluyen deliberadamente las variables de reloj (hora, minuto, segundo) para evitar que el
modelo aprenda atajos temporales triviales.

\chapter{Hiperparámetros de los modelos}
\label{anexo:hiperparametros}

\begin{table}[h]
\centering
\caption{Hiperparámetros principales del especialista (HistGradientBoosting, dos cabezas).}
\begin{tabular}{ll}
\toprule
\textbf{Hiperparámetro} & \textbf{Valor} \\
\midrule
Tasa de aprendizaje & 0,045 \\
Iteraciones máximas & 220 (con parada anticipada) \\
Fracción de validación interna & 0,15 \\
Máximo de hojas por árbol & 15 \\
Mínimo de muestras por hoja & 120 \\
Regularización L2 & 0,10 \\
\bottomrule
\end{tabular}
\end{table}

\begin{table}[h]
\centering
\caption{Hiperparámetros principales de los modelos profundos (línea de exploración).}
\begin{tabular}{ll}
\toprule
\textbf{Hiperparámetro} & \textbf{Valor} \\
\midrule
Optimizador & AdamW \\
Tasa de aprendizaje & 0,0015 \\
Decaimiento de pesos & $10^{-4}$ \\
Tamaño de lote & 192 \\
Paciencia (parada anticipada) & 18 épocas \\
Dimensión oculta de la GRU & 56 \\
\emph{Dropout} (\emph{BookEncoder}) & 0,08 \\
\bottomrule
\end{tabular}
\end{table}

La política congelada opera con dos umbrales: valor esperado predicho superior a 0,75 y
probabilidad de salud no inferior a 0,50. El filtro de volatilidad habilita la operativa cuando
la volatilidad realizada a 5 segundos no supera 0,6657 (percentil 80 de la volatilidad de
entrenamiento). El intervalo de confianza se estima por \emph{bootstrap} con 5\,000 remuestreos
y semilla fija.
